\scanpage{403}\mspage{317}

il faut construire

1) $i_\alpha\in\mathrm{Ob}\,I$ et un hom $(i_\beta)_{\beta\in I_\alpha}\longrightarrow i_\alpha$ i.e. des $\tau_{\deleted{\unsure{\beta'}}\alpha,\beta}:i_\beta\unsure{-|\longrightarrow}i_\alpha$, \inserted{$(\beta\in I_\alpha)$}, compatibles aux transitions $\tau_{\beta',\beta}$ pour $\beta<\beta'<\alpha$

2) $j_\alpha\in\mathrm{Ob}\,J'$, et un hom $(j_\beta)_{\beta\in I_\alpha}\unsure{-|\longrightarrow}j_\alpha$, i.e des $\tau'_{\alpha,\beta}:i_\beta\unsure{-|\longrightarrow}i_\alpha$, $(\beta\in I_\alpha)$, compatibles aux transitions $\tau'_{\beta',\beta}$ pour $\beta<\beta'<\alpha$ --

3) \deleted{\unsure{Une}} flèche

\[
u_\alpha:\unsure{\varphi|\phi}(i_\alpha)\unsure{-|\longrightarrow}j_\alpha
\]

4) Si $\alpha$ n'est pas ordinal limite, i.e. $\alpha=\unsure{\alpha'_0|\alpha_0}+1$, une flèche

\[
v_{\alpha_0}:\deleted{\unsure{\varphi}}j_{\alpha_0}\longrightarrow\varphi(i_{\alpha_0+1})=\varphi(i_\alpha)\,;
\]

\deleted{de façon} et ceci de façon à satisfaire aux contraintes \textcircled{1}\textcircled{2}\textcircled{3}\textcircled{4} quand on y suppose $\beta+1\leq\alpha$ (ald $\beta+1<\alpha$) dans \textcircled{1} et \textcircled{2}, et $\beta'\leq\alpha$ (ald $\beta'<\alpha$) dans \textcircled{3}, et $\beta'+1\leq\alpha$ (ald $\beta'+1<\alpha$) dans \textcircled{4}. Donc \inserted{pour} les contraintes $1^\circ,2^\circ,4^\circ$, \deleted{\unsure{vient à être}\ \unsure{en}} il n'y a en tenir compte \uline{que} si $\alpha$ n'est pas ordinal limite, donc $\alpha=\alpha_0+1$, et en y faisant

\[
\beta'=\alpha_0.
\]

Quant à la contrainte $3^\circ$, il faut en

% archive sheet number: 402
